% appendices/appendix-d-sumber-belajar.tex
\chapter{Sumber Belajar}
\label{app:sumber}

Lampiran ini menjelaskan kegunaan setiap rujukan pada Daftar Pustaka, agar
Anda dapat memilih bacaan menurut keperluan. Uraian lengkap setiap rujukan
tidak diulang di sini; yang ditambahkan hanyalah keterangan tentang bagian
mana yang berguna dan untuk bab mana.

\section{Buku Teks dan Dokumentasi Resmi}

Rujukan pada Tabel~\ref{tab:sumber-utama} adalah bacaan pokok. Ketiganya
tersedia bebas, dan ketiganya dipakai berulang kali sepanjang buku ini.

\begin{table}[htbp]
\centering
\caption{Buku teks dan dokumentasi resmi beserta kegunaannya}
\label{tab:sumber-utama}
{\setlength{\tabcolsep}{4pt}\footnotesize%
\begin{tabularx}{\textwidth}{@{}C{1.1cm} L{4.6cm} Y@{}}
\toprule
\headrow \textbf{Kode} & \textbf{Bacaan} & \textbf{Kegunaan pada buku ini} \\
\midrule
\cite{sweigart-mg} & Sweigart, \istilah{Making Games with Python \& Pygame} &
Pembahasan permainan lengkap dengan Pygame, dari \istilah{game loop} sampai
bunyi. Berguna sebagai contoh utuh setelah Bab~XI. \\
\addlinespace
\cite{sweigart-invent} & Sweigart, \istilah{Invent Your Own Computer Games with
Python} & Dasar Python dan permainan berbasis teks. Berguna sebelum Bab~III
apabila dasar Python Anda belum kuat. \\
\addlinespace
\cite{pygame-docs} & Dokumentasi resmi Pygame &
Rujukan perintah yang dipakai pada Bab~III--XI. Periksa di sini ketika nama
atau argumen sebuah fungsi tidak lagi sesuai dengan contoh pada buku. \\
\addlinespace
\cite{pygame-sprite-mask} & Dokumentasi \code{pygame.mask} &
Pembahasan tabrakan berdasarkan bentuk gambar, yang tidak dibahas pada
Bab~VII tetapi berguna untuk permainan dengan bentuk tidak beraturan. \\
\addlinespace
\cite{unity-leveldesign} & Unity Technologies, \istilah{Introduction to Level
Design in Game Development and in Unity} &
Rancangan level, penyampaian keterangan kepada pemain, dan pengujian. Menjadi
sandaran Bab~XII--XVI. \\
\addlinespace
\cite{unity-learn} & Unity Learn &
Kumpulan tutorial resmi Unity. Berguna untuk melatih ulang langkah pada
Bab~XII--XV sebelum mengerjakan proyek. \\
\addlinespace
\cite{unity-manual-import} & Unity Manual: \istilah{Importing Assets} dan
\istilah{Audio} &
Rujukan pengaturan impor dan bunyi pada Bab~XIV--XV. \\
\bottomrule
\end{tabularx}}
\end{table}

\section{Tulisan Teknis Mendalam}

Rujukan pada Tabel~\ref{tab:sumber-teknis} membahas satu topik secara
mendalam. Sebagian besar tidak perlu dibaca seluruhnya; bacalah bagian yang
bersangkutan ketika Anda ingin memahami alasan di balik sebuah pilihan.

\begin{table}[htbp]
\centering
\caption{Tulisan teknis mendalam beserta bab yang berkaitan}
\label{tab:sumber-teknis}
{\setlength{\tabcolsep}{4pt}\footnotesize%
\begin{tabularx}{\textwidth}{@{}C{1.1cm} L{4.6cm} Y@{}}
\toprule
\headrow \textbf{Kode} & \textbf{Bacaan} & \textbf{Kegunaan pada buku ini} \\
\midrule
\cite{nystrom-gameloop} & Nystrom, ``Fix Your Timestep'' dan pembahasan
\istilah{game loop} & Alasan di balik \istilah{game loop} pada Bab~IV, termasuk
pemisahan antara pembaruan tetap dan pembaruan tampilan. \\
\addlinespace
\cite{fiedler-timestep} & Fiedler, \istilah{Timestep} &
Perbandingan beberapa cara menjalankan waktu pada permainan. Berguna setelah
Bab~IV apabila gerak permainan Anda tidak stabil. \\
\addlinespace
\cite{catto-integration} & Catto, \istilah{Integration Basics} &
Penjelasan Euler dan Verlet untuk fisika sederhana. Sandaran Bab~VI, termasuk
alasannya mengenai galat yang mengecil ketika langkah waktunya diperkecil. \\
\addlinespace
\cite{santos-pathfinding} & Santos, \istilah{Pathfinding} &
Penjelasan BFS, Dijkstra, dan A* beserta ilustrasinya. Sandaran Bab~X. \\
\addlinespace
\cite{santos-fsm} & Santos, \istilah{Finite State Machines} &
Penjelasan mesin keadaan beserta contohnya. Sandaran Bab~IX. \\
\addlinespace
\cite{bendersky-pygame-astar} & Bendersky, A* pada Pygame &
Contoh penerapan A* pada permainan dua dimensi. Berguna sebagai pembanding
\code{ch10_astar.py}. \\
\addlinespace
\cite{munir-astar} & Munir, \istilah{Pencarian Jalur dengan A*} &
Pembahasan sifat heuristik dan pengaruhnya pada hasil pencarian. Berguna untuk
menjawab latihan pada Bab~X. \\
\addlinespace
\cite{itb-astar-navigasi} & ITB, ``Implementasi Algoritma A* dalam Navigasi
Karakter'' & Contoh penerapan A* pada karakter permainan. \\
\addlinespace
\cite{stmikdharmapala-astar} & STMIK Dharmapala, penerapan A* pada permainan &
Contoh penerapan pada permainan buatan Indonesia. \\
\addlinespace
\cite{itb-collision} & ITB, ``Collision Detection in Precision Parry-Based
Combat'' & Pembahasan deteksi tabrakan yang lebih rinci daripada Bab~VII. \\
\addlinespace
\cite{pygame-chimp} & Pygame Community, \istilah{Chimp Line-by-Line} &
Contoh permainan Pygame yang dibahas baris demi baris. Berguna setelah
Bab~VII untuk melihat susunan program secara utuh. \\
\addlinespace
\cite{reintech-game-ai} & Reintech, ``How to Create a Video Game AI with
Python'' & Contoh perilaku NPC dalam Python. Berguna sebagai pembanding
Bab~VIII. \\
\bottomrule
\end{tabularx}}
\end{table}

\section{Rencana Pembelajaran dan Kurikulum Pembanding}

Rujukan pada Tabel~\ref{tab:sumber-rps} dipakai untuk menyusun peta capaian
pada Bab~XVII, dan berguna bagi Anda yang ingin mengetahui bagaimana mata
kuliah serupa disusun di tempat lain.

\begin{table}[htbp]
\centering
\caption{Rencana pembelajaran dan kurikulum pembanding}
\label{tab:sumber-rps}
{\setlength{\tabcolsep}{4pt}\footnotesize%
\begin{tabularx}{\textwidth}{@{}C{1.1cm} L{4.6cm} Y@{}}
\toprule
\headrow \textbf{Kode} & \textbf{Bacaan} & \textbf{Kegunaan pada buku ini} \\
\midrule
\cite{its-rps-game} & ITS, \istilah{RPS Desain dan Pemrograman Game} &
Pembanding susunan capaian dan pokok bahasan. \\
\addlinespace
\cite{unima-rps-game} & UNIMA, \istilah{RPS Pengembangan Game} &
Pembanding bentuk penilaian dan presentasi proyek. \\
\addlinespace
\cite{unesa-rps-game} & UNESA, \istilah{RPS Pemrograman Game} &
Pembanding urutan materi antara blok algoritma dan blok mesin permainan. \\
\addlinespace
\cite{polibatam-kurikulum} & Politeknik Negeri Batam, \istilah{Kurikulum
Program Studi Teknologi Permainan} & Pembanding daftar periksa capaian dan
penyusunan daftar periksa pada Bab~XVI. \\
\addlinespace
\cite{univrab-pemrograman-game} & Universitas Abdurrab, \istilah{Deskripsi Mata
Kuliah Pemrograman Game} & Pembanding cakupan materi dan prasyarat. \\
\bottomrule
\end{tabularx}}
\end{table}

\section{Karya Akademik Indonesia}

Rujukan pada Tabel~\ref{tab:sumber-karya} adalah skripsi, jurnal, dan laporan
pengabdian dari perguruan tinggi Indonesia. Kegunaannya berbeda dari buku
teks: karya seperti ini memperlihatkan bagaimana sebuah permainan sungguh
dibangun dan dilaporkan, termasuk kekurangannya.

\begin{table}[htbp]
\centering
\caption{Karya akademik Indonesia beserta kegunaannya}
\label{tab:sumber-karya}
{\setlength{\tabcolsep}{4pt}\footnotesize%
\begin{tabularx}{\textwidth}{@{}C{1.1cm} L{4.6cm} Y@{}}
\toprule
\headrow \textbf{Kode} & \textbf{Bacaan} & \textbf{Kegunaan pada buku ini} \\
\midrule
\cite{uma-pygame} & Universitas Medan Area, \istilah{Panduan Membuat Game 2D
Sederhana dengan Pygame} & Contoh pengantar yang ringkas untuk Bab~III--IV. \\
\addlinespace
\cite{ub-pygame} & Universitas Brawijaya, ``Game Development with Python'' &
Contoh penyusunan permainan sederhana dan penjelasannya. \\
\addlinespace
\cite{dinus-pygame} & Universitas Dian Nuswantoro, ``Pemrograman dan Game
dengan Python/Pygame'' & Contoh pelatihan Pygame bagi pemula. \\
\addlinespace
\cite{stmikjayakarta-pygame} & STMIK Jayakarta, ``Penerapan Library Pygame
dalam Game RPG The Adventure'' & Contoh permainan RPG dengan Pygame, termasuk
pengurusan bunyi. \\
\addlinespace
\cite{itn-astar-fsm} & ITN Malang, ``Penerapan A* dan FSM pada Game Lost
Civilization'' & Contoh penggabungan A* dan FSM, sebagaimana Bab~XI. \\
\addlinespace
\cite{umri-astar-behaviour} & UMRI, ``Penerapan Algoritma A* dan Behaviour
Trees'' & Pembanding pendekatan lain untuk perilaku NPC. \\
\addlinespace
\cite{uinmalang-fsm-fuzzy} & UIN Malang, \istilah{Implementasi FSM dan Logika
Fuzzy} & Pembanding pendekatan yang lebih halus daripada FSM pada Bab~IX. \\
\addlinespace
\cite{uny-katelu} & UNY, \istilah{Pengembangan Permainan Edukasi KATELU dengan
Unity 3D} & Contoh proyek Unity lengkap beserta laporannya. \\
\addlinespace
\cite{undip-media} & UNDIP, ``Media Pembelajaran Budaya Indonesia Menggunakan
Unity'' & Contoh pencantuman aset dan sumber pada karya Unity. \\
\addlinespace
\cite{untan-edukasi} & UNTAN, ``Game Edukasi Budaya Indonesia'' &
Pembanding rancangan permainan edukasi. \\
\addlinespace
\cite{rokania-unity} & STKIP Rokania, \istilah{Pelatihan Dasar Pemrograman Game
dengan Unity} & Contoh pelatihan Unity bagi pemula, dan pembanding refleksi
pemilihan teknologi pada Bab~XVI. \\
\bottomrule
\end{tabularx}}
\end{table}

\section{Saran Urutan Membaca}

Buku ini disusun agar dapat dibaca berurutan. Apabila Anda memerlukan bacaan
tambahan, urutan berikut biasanya paling berguna.

\begin{enumerate}
  \item \textbf{Sebelum Bab~III}: \cite{sweigart-invent} untuk dasar Python.
  \item \textbf{Setelah Bab~IV}: \cite{nystrom-gameloop} dan
        \cite{fiedler-timestep} untuk memahami mengapa waktunya dijalankan
        dengan cara itu.
  \item \textbf{Setelah Bab~VI}: \cite{catto-integration} untuk galat
        integrasi dan cara memperkecilnya.
  \item \textbf{Setelah Bab~IX dan Bab~X}: \cite{santos-fsm} dan
        \cite{santos-pathfinding} untuk pembahasan yang lebih rinci.
  \item \textbf{Setelah Bab~XI}: \cite{sweigart-mg} untuk melihat satu
        permainan Pygame yang lengkap.
  \item \textbf{Setelah Bab~XVI}: \cite{unity-leveldesign} bagian pengujian,
        lalu salah satu karya akademik pada Tabel~\ref{tab:sumber-karya}
        sebagai contoh proyek yang dilaporkan.
\end{enumerate}
